% Generated draft: requires mathematical and source review.
\documentclass[11pt]{article}
\usepackage[margin=25mm]{geometry}
\usepackage{fontspec}
\setmainfont{Latin Modern Roman}
\usepackage{amsmath,amssymb,mathtools,unicode-math}
\setmathfont{Latin Modern Math}
\usepackage{enumitem,needspace,xurl,hyperref}
\setlength{\parindent}{0pt}
\setlength{\parskip}{5pt}
\setlength{\emergencystretch}{3em}
\newcommand{\N}{\mathbb{N}}
\newcommand{\Z}{\mathbb{Z}}
\newcommand{\Q}{\mathbb{Q}}
\newcommand{\R}{\mathbb{R}}
\newcommand{\C}{\mathbb{C}}
\newcommand{\F}{\mathbb{F}}
\newcommand{\A}{\mathbb{A}}
\newcommand{\PP}{\mathbb{P}}
\newcommand{\E}{\mathbb{E}}
\newcommand{\eps}{\varepsilon}
\newcommand{\dd}{\,\mathrm{d}}
\begin{document}
\section*{UnsolvedMath ID: 30006138}
\textbf{Source problem code:} catalogue entry
\subsection{Definitions and mathematical statement}
Let \(n\geq 1\) be an integer. A convex body is a compact convex set \(K\subset\R^n\) with nonempty interior. Writing \(|K|\) for its \(n\)-dimensional Lebesgue measure, define its centroid by
\[
c(K)=\frac{1}{|K|}\int_K x\,\dd x.
\]
For an affine hyperplane \(H\subset\R^n\) containing \(c(K)\), define the centroid of the section using the induced \((n-1)\)-dimensional measure:
\[
c_H(K\cap H)=\frac{\displaystyle\int_{K\cap H}x\,\dd\mathcal H^{n-1}(x)}{\mathcal H^{n-1}(K\cap H)}.
\]
Here \(\mathcal H^{n-1}\) denotes Hausdorff measure with its usual Euclidean normalization; when \(n=1\), \(\mathcal H^0\) is counting measure. Since \(c(K)\) lies in the interior of \(K\), the denominator is positive. Set
\[
\mathcal S(K)=\bigl\{H\subset\R^n: H\text{ is an affine hyperplane},\ c(K)\in H,\ c_H(K\cap H)=c(K)\bigr\},
\qquad
\mu(n)=\min_{\substack{K\subset\R^n\\K\text{ a convex body}}}\#\mathcal S(K).
\]
Distinct sections are counted by distinct affine hyperplanes, without orientation; no finiteness of \(\mathcal S(K)\) is assumed.

The original question asks whether, for every convex body \(K\subset\R^n\),
\[
\#\mathcal S(K)\geq n+1,
\]
and, more generally, asks for the value of \(\mu(n)\).

The supplied literature assessment, dated \(2026\text{-}08\text{-}22\), reports
\[
\mu(2)=3,\qquad \mu(n)=1\quad(n\geq5),
\]
and identifies \(\mu(3)\) and \(\mu(4)\) as unresolved. These are supplied dated curation claims, not independently verified mathematical results in this draft. The universal \(n+1\) question is retained as the original question; the supplied assessment would rule out that bound in dimensions \(n\geq5\). The record does not explicitly state a lower bound on the dimension; the definitions above use the standard positive-integer convention, including the zero-dimensional section convention when \(n=1\).
\subsection{Short English statement}
Count the hyperplane sections through a convex body's centroid whose own centroids equal that point, and determine the smallest possible count in each dimension. The supplied dated assessment identifies dimensions \(3\) and \(4\) as the remaining open cases; its literature claims require independent review.
\subsection{Sources}
\begin{itemize}
\item Supplied local UnsolvedMath registry record, permanent ID \(30006138\), title: Multiplicity of Centroidal Hyperplane Sections; exact source problem code: \texttt{catalogue entry}. Collection: Oberwolfach Reports Open Problems. The record gives proposed year \(2024\) and names no proposer.
\item Supplied source citation: \emph{Convex Geometry and its Applications} (2025), Oberwolfach Reports, pp.\ 3301--3376, DOI \texttt{10.4171/owr/2024/58}. \url{https://doi.org/10.4171/owr/2024/58}. The citation's year \(2025\) and the record's proposed year \(2024\) are preserved as supplied.
\item Supplied literature reference: Myroshnychenko--Tatarko--Yaskin, \emph{Adv. Math.} 2025. \url{https://doi.org/10.1016/j.aim.2024.110081}.
\item Supplied literature reference: Haddad--Jimenez--Villa, 2026. \url{https://doi.org/10.1007/s12220-026-02570-2}.
\item Literature assessment date: \texttt{2026-08-22}. The supplied dated curation labels the problem partially solved and reports \(\mu(2)=3\), \(\mu(n)=1\) for \(n\geq5\), and unresolved values in dimensions \(3\) and \(4\). This assessment is not independently verified here. The background calls this the Loewner-centroid minimum, without defining that terminology separately.
\item The registry labels the statement \texttt{corrected\_verified} and says it was reconstructed and checked against the official Oberwolfach report. This is supplied provenance, not a source-verification claim by this draft.
\end{itemize}
\end{document}
